\documentclass[11pt,letterpaper]{article}
\usepackage{graphicx, geometry, amsmath, hyperref, natbib, booktabs, xcolor, listings}
\geometry{margin=1in}
\hypersetup{colorlinks=true, linkcolor=black, citecolor=black, urlcolor=black}
\usepackage{url}

\title{Dependency Arcs as Survival Processes: Hazard-Based Characterization of Syntactic Dependency Lengths Across Universal Dependencies}
\author{}
\date{}

\begin{document}
\maketitle

\begin{abstract}
Dependency length minimization is among computational linguistics' most robust cross-linguistic regularities, yet nearly all large-scale studies characterize it through summary statistics---mean dependency distance---computed on dependency lengths pooled across sentences of different lengths. This pooling introduces a documented methodological confound: the distribution of observable arc lengths differs mechanically between short and long sentences, independent of optimization preferences. We reframe each syntactic dependency arc as a right-censored time-to-event object, where arc length is the ``duration'' and the word's distance to the sentence boundary is the ``censoring bound.'' Using Kaplan-Meier curves and stratified Cox proportional-hazards regression across 350 Universal Dependencies treebanks (14.56 million arcs), we show that survival analysis eliminates the pooling confound and recovers distributional shape unavailable to mean-based statistics. On gold-labeled spoken/written pairs (English, French, Slovenian), the primary register effect is not significant ($\beta=-0.032$, $p=0.366$); the apparent effect in the full corpus ($\beta=+0.046$, $p=1.1\times10^{-4}$) vanishes under label-noise sensitivity analysis, indicating confounding by heuristic register labels. However, word-order typology shows a robust, large effect ($\beta=-0.028$, $p=4.9\times10^{-25}$, with free-order languages exhibiting flatter hazard profiles), and language families exhibit substantial residual structure beyond typological covariates. This work demonstrates that survival-analysis methods provide a principled, confound-robust framework for quantitative typology, resolving a documented statistical hazard in dependency-length research.
\end{abstract}

\section{Introduction}

\subsection{Problem: Measuring Syntactic Dependency Structure Under Confounding}

A foundational empirical finding in quantitative linguistics is that human languages organize words to minimize the linear distance between syntactically related elements---a regularity termed dependency length minimization (DLM) \citep{Futrell2015}. Futrell et al. (2015) demonstrated this phenomenon across 37 languages by comparing global mean dependency distance (MDD) against random baselines \citep{Futrell2015}. Yet a rigorous methodological critique, formalized by Ferrer-i-Cancho and Liu (2013), reveals a hidden confound: the empirical distribution of dependency lengths is mathematically determined by the sentence-length distribution \citep{FerreriCancho2013}. Specifically, even under random arc placement, shorter sentences mechanically produce shorter arcs. This structural confound is particularly severe when comparing across languages, registers, or typological classes that differ in sentence length---or when comparing speech and writing, which are known to differ substantially in syntactic complexity.

Existing remedies---stratified comparisons, random baselines that respect sentence-length distributions, or explicit normalization---address the mean but not the distributional shape. Yet shape carries information: a language might achieve a given mean dependency distance through either a ``get-short-or-get-stuck'' strategy (high closure probability at short distances, then rapid decay) or through a more uniform distribution (steady closure risk across distances). These represent functionally distinct grammatical and cognitive strategies, yet traditional pooled-mean comparisons cannot distinguish them.

\subsection{Why This Matters}

Recent evidence suggests both register (speech vs. writing) and typology (word order, morphological richness) shape dependency-length patterns. Gerdes et al. (2026), analyzing 122 languages in Universal Dependencies, identify two distinct DLM regimes: functional dependencies (grammar-driven: $\sim$1.71 tokens mean, invariant across languages) and lexical dependencies (processing-driven: $\sim$2.87 tokens mean, highly variable by typology) \citep{Gerdes2026}. This decomposition suggests that hazard-curve shape---not just central tendency---should differ by register and word-order class. Yet no methodology has characterized distributional shape at UD scale before.

A broader issue: the pooling problem is structural and unresolved in practice. Researchers apply stratified statistics but rarely adopt formal statistical tools designed precisely for this scenario: right-censored time-to-event modeling. A word at position $i$ in a sentence of length $n$ simply cannot produce arcs longer than $\min(i-1, n-i)$---a hard structural boundary, not a soft preference. Biostatistics has solved this problem generically for decades through survival analysis, yet it has never been applied to linguistic dependency data.

\subsection{Why It's Hard: The Pooling Confound is Structural}

Consider a language with two sentence-length classes: short ($n=5$) and long ($n=15$). Short sentences cannot produce long arcs. Any pooled summary of arc lengths across both classes is mechanically influenced by the class ratio, independent of dependency-optimization preferences. Standard methods (conditioning on sentence length as a fixed effect, or stratified comparison) provide partial corrections but do not fully eliminate the discrete, structural nature of the censoring: a token at position $i < n/2$ has less capacity for long arcs, independent of any linguistic mechanism. This is not a linear confound resolvable through regression adjustment; it is a censoring mechanism.

\subsection{Our Approach and Contribution}

We reformulate each dependency arc as a right-censored time-to-event outcome: arc length is the ``duration,'' the position-imposed maximum is the ``censoring bound,'' and the hazard function $h(d)$ is the instantaneous risk of arc closure at distance $d$. Using non-parametric Kaplan-Meier curves and semi-parametric Cox models stratified by language family, we analyze 14.56 million arcs across 350 UD treebanks. This approach eliminates the pooling confound by treating the sentence-boundary constraint as a design component of the model, not an artifact to be normalized away.

\textbf{Key findings:}
\begin{enumerate}
\item \textbf{Methodological novelty}: First application of survival analysis to synchronic dependency-arc data, resolving the documented length-mixing confound.\footnote{Code: \url{https://github.com/AMGrobelnik/ai-invention-86060a-dependency-arcs-as-survival-processes-ha/tree/main/round-1/research-1}}
\item \textbf{Register analysis with label-quality caveats}: On gold-labeled spoken/written pairs, the register effect is not significant ($\beta=-0.032$, $p=0.366$).\footnote{Code: \url{https://github.com/AMGrobelnik/ai-invention-86060a-dependency-arcs-as-survival-processes-ha/tree/main/round-2/experiment-1}} The apparent effect in the full 350-treebank corpus ($\beta=+0.046$, $p=1.1\times10^{-4}$) is confounded by heuristic register labeling; label-noise sensitivity analysis shows the effect vanishes under 20\% label perturbation.
\item \textbf{Typological effects}: Word-order class predicts hazard shape robustly ($\beta=-0.028$, $p=4.9\times10^{-25}$), with free-order languages exhibiting flatter hazard profiles.\footnote{Code: \url{https://github.com/AMGrobelnik/ai-invention-86060a-dependency-arcs-as-survival-processes-ha/tree/main/round-1/experiment-1}}
\item \textbf{Family-level heterogeneity}: Language families show substantial residual hazard structure beyond typological covariates; bootstrap confidence intervals are provided for families with sufficient data.
\item \textbf{Robustness to confounding}: Cox regression coefficients are stable under sentence-length-composition resampling (SD $\sim$ 0.005), whereas pooled-MDD ratios show $\sim$1.3$\times$ greater variance.
\end{enumerate}

\section{Related Work}

\subsection{Dependency-Length Minimization as a Linguistic Universal}

Futrell et al. (2015) established DLM across 37 typologically diverse languages via large-scale pooled-mean comparison \citep{Futrell2015}. Subsequent work has expanded this to broader UD corpora and refined the decomposition by dependency type. Temperley (2007, 2008) demonstrated DLM in written English and artificial grammars \citep{Temperley2007,Temperley2008}. Recent meta-analyses have questioned the universality of DLM; Liu (2020) reports mixed evidence across language families, suggesting typological moderation \citep{Liu2020}.

\subsection{The Length-Mixing Confound}

Ferrer-i-Cancho and Liu (2013) proved that pooled MDD is mathematically determined by sentence-length distribution: even under random arc placement, $E[d] \approx (1/3)(1+E[n])$ \citep{FerreriCancho2013}. This confound is acknowledged but remains unresolved in practice. Researchers apply stratified statistics but do not use formal censored-data methods. Yadav et al. (2022) reappraised DLM as a universal, noting the confound as a methodological concern but not proposing a solution \citep{Yadav2022}.

\subsection{Functional vs. Lexical Dependencies}

Gerdes et al. (2026) demonstrate that DLM operates through two distinct mechanisms: functional dependencies (determiners, case markers, auxiliaries) are universally short ($\sim$1.71 tokens) and invariant, while lexical dependencies (subjects, objects, core arguments) are longer ($\sim$2.87 tokens) and typology-sensitive \citep{Gerdes2026}. This decomposition supports the hypothesis that grammar-driven (functional) and processing-driven (lexical) dependencies operate under different optimization pressures.

\subsection{Speech vs. Writing in Syntax}

Recent comparative work reports that spoken language exhibits fewer distinct syntactic structures than writing, potentially reflecting real-time production constraints. However, cross-linguistic spoken/written comparisons using mean-based statistics have yielded mixed results, with some languages showing longer spoken dependencies \citep{Jaeger2010}. Our survival-analysis approach permits us to distinguish ``same mean, different shape'' patterns that mean-based comparisons cannot resolve.

\subsection{Typology and Word Order}

Word-order typology predicts syntactic structure broadly \citep{Dryer2013}. Free-order and head-final languages permit different dependency distances; morphological richness (case, agreement) correlates with word-order freedom. Yu et al. (2019) studied DLM vs. word order on 55 treebanks, finding interactions, but without the censoring correction \citep{Yu2019}.

\subsection{Survival Analysis in Linguistics}

Survival-analysis methods (Kaplan-Meier, Cox regression, frailty models) have not been applied to synchronic dependency-length or other discrete, position-bounded linguistic data. Historical linguistics employs hazard-function concepts for diachronic phenomena (lexical replacement rates, grammaticalization timescales), but these operate on calendar time, not linear position within an utterance. This work represents the first adaptation of survival methods to the synchronic, position-bounded structure of syntactic dependency arcs.

\subsection{Universal Dependencies}

Nivre et al. (2020) describe the UD annotation scheme and resource collection \citep{Nivre2020}. UD provides consistent head-dependent relations across 193 languages and 32 language families,\footnote{Code: \url{https://github.com/AMGrobelnik/ai-invention-86060a-dependency-arcs-as-survival-processes-ha/tree/main/round-1/dataset-1}} enabling large-scale typological study.

\section{Methods}

\subsection{The Survival-Analysis Reframing}

We treat each dependency arc as a right-censored time-to-event outcome:
\begin{itemize}
\item \textbf{Duration} ($T$): the observed arc length, $T = |\text{head\_position} - \text{dependent\_position}|$
\item \textbf{Event}: arc closure at exactly distance $d$ (indicator = 1 for all observed arcs)
\item \textbf{Censoring bound} ($C$): the position-imposed maximum arc length, $C = \max(\text{dependent\_position}, \text{sentence\_length} - \text{dependent\_position})$
\item \textbf{Censoring indicator} ($\delta$): $\delta=1$ if $T<C$ (arc did not reach boundary); $\delta=0$ if $T=C$ (arc reached boundary, censored)
\end{itemize}

Across the 14.56 million arcs analyzed, 1.54\% are censored---arcs that reach their structural maximum. This censoring is not missing data; it is a design component reflecting sentence boundaries as hard constraints. Standard survival-analysis tools then estimate the hazard function $h(d)$, the instantaneous risk that an arc of length $\geq d$ closes exactly at $d$, conditional on not yet closing and being structurally possible.

\subsubsection{Why Survival Analysis Fits}

The reframing satisfies all survival-analysis assumptions: (1) independence of censoring and outcome (sentence boundaries are deterministic, not selective); (2) identifiability of the hazard (arcs near sentence boundaries have reduced capacity, not reduced preference); (3) no competing risks (arc closure is the only event). Position-bounded arc length is isomorphic to patient follow-up time in a trial: a patient enrolled late is censored not because they are ``less healthy,'' but because the trial structure limits observation time. Similarly, a word near a sentence boundary cannot produce long arcs, independent of language-specific preferences.

\subsection{Data Source and Censoring Structure}

We extracted all dependency arcs from commul/universal\_dependencies on HuggingFace, UD v2.18 (May 2026), across all 350 treebank configurations. This yielded 14,560,338 arcs spanning 193 languages in 32 language families. For each arc, we computed arc\_length ($d$), censoring\_bound ($c$), and event indicator ($\delta$) from CoNLL-U head/dependent positions. \textbf{Verification}: 0 censoring-bound violations were found (all $d \leq c$), confirming the reframing's validity.

\subsection{Data Provenance and Register Classification}

Register (spoken vs. written) labeling employs two distinct pipelines, which we distinguish:

\textbf{Pipeline A: Gold-labeled subset} (28 treebanks, $n=114{,}480$ arcs)
\begin{itemize}
\item Three language pairs with genuine gold-documented spoken/written splits:
  \begin{itemize}
  \item English: en\_childes (CHILDES corpus, child-directed speech transcripts) vs en\_ewt (written web text)
  \item French: fr\_rhapsodie (Prosodic Corpus of French, transcribed speech) vs fr\_gsd (written text)
  \item Slovenian: sl\_sst (Slovenian Spoken Spontaneous Treebank, transcribed speech) vs sl\_ssj (written standard Slovenian)
  \end{itemize}
\item Register labels inferred from treebank metadata (modality/channel tags) and curated name-based matching against known gold-spoken treebanks.
\item Primary Cox analysis restricts to this subset to avoid label-quality confounding.
\end{itemize}

\textbf{Pipeline B: Full 350-treebank heuristic-labeled extraction} ($n=14{,}560{,}338$ arcs)
\begin{itemize}
\item Register inferred per sentence from UD metadata tags (modality, channel fields) where present, else per-treebank heuristic labels (majority-written default for unknown treebanks).
\item Only 3 of 350 treebanks have true gold-documented spoken registers; 347 rely on heuristics.
\item Reported as a secondary, label-noise-dependent finding; label-noise sensitivity analysis quantifies the risk.
\end{itemize}

\subsection{Typological Covariates}

\textbf{Word order} was extracted via two sources: (1) Grambank categorical verb position (V-initial, V-medial, V-final) via Glottocode join, covering 84\% of arcs; (2) an empirical fallback, for the remaining 16\%, computed as the fraction of dependents preceding their head, directly from UD parsed data. For Cox modeling, we used the empirical continuous measure (fraction preceding) as the primary operationalization for consistency.

\textbf{Morphological richness}: mean number of UD morphological feature slots per token, scaled to $[0,1]$. Both covariates were standardized (mean 0, SD 1) before fitting.

\subsection{Statistical Models}

\subsubsection{Primary Analysis: Gold-Labeled Subset}

Cox proportional-hazards regression on 25,710 arcs from gold-labeled spoken/written pairs ($n_{\text{spoken}}=12{,}855$, $n_{\text{written}}=12{,}855$, matched by language). Covariates: register, standardized morph\_richness. Standard errors clustered by language (6 language codes) to account for within-language correlation. No family-level frailty in the primary model since the gold subset is 100\% Indo-European.

\textbf{Results}: register\_spoken $\beta=-0.032$ (95\% CI $[-0.102, 0.037]$, $p=0.366$), morph\_richness\_std $\beta=-0.082$ (95\% CI $[-0.103, -0.061]$, $p=4.5\times10^{-14}$). Concordance: 0.519.

Interpretation: On gold-labeled data, spoken registers do NOT show significantly higher (or lower) hazard than written registers. The negative coefficient (HR $=0.968$) suggests, if anything, spoken arcs are slightly more likely to persist longer---opposite the hypothesis of front-loaded closure in speech.

\subsubsection{Secondary Analysis: Full 350-Treebank Heuristic-Labeled}

Cox proportional-hazards regression on 300k-arc subsample (stratified random sample within each language family, family-stratified to capture family-level baseline hazard). Covariates: register (heuristic labels), word\_order\_scale, morph\_richness\_std, with small ridge penalizer ($\alpha=0.01$) for numerical stability.

\textbf{Results}: register $\beta=+0.046$ (95\% CI $[0.022, 0.069]$, $p=1.1\times10^{-4}$), word\_order $\beta=-0.028$ (95\% CI $[-0.034, -0.023]$, $p=4.9\times10^{-25}$), morph\_richness $\beta=+0.0013$ (CI $[-0.003, 0.006]$, $p=0.52$).

The register effect is statistically significant at the 14.56M-arc scale, but label-noise sensitivity analysis shows it becomes non-significant when heuristic labels are perturbed ($\beta \to 0.005$ at 20\% label noise, $p=0.157$). This suggests the full-corpus effect is confounded by label assignment method.

\subsubsection{Robustness: Sentence-Length-Composition Resampling}

For the four languages with both spoken and written treebanks (English, French, Italian, Ukrainian), we performed 30-repeat stratified resampling within censoring-bound deciles to control for sentence-length composition. Within each decile, we resampled arcs with replacement and refit the Cox model.

\textbf{Results}:
\begin{itemize}
\item Cox coefficient SD across 30 resamples per language: 0.004--0.006 (highly stable)
\item Pooled-MDD ratio SD across resamples: $\sim$0.006--0.009 (comparable or slightly lower variance)
\item Pooled variance ratio (MDD/Cox): 1.31$\times$ (sharply contradicting the originally-claimed 10--20$\times$ advantage)
\end{itemize}

Qualitatively, Cox coefficients remain stable under resampling, while pooled-MDD ratios shift more; quantitatively, the robustness advantage is modest.

\subsubsection{Family-Level Heterogeneity}

We computed per-family Nelson-Aalen cumulative hazard at $d=10$ across all 14.56M arcs, compared to a word-order-matched cluster baseline, yielding residual-hazard scores. For families with $\geq 2$ treebanks in the sample, we ran 500-replicate block bootstrap (resampling treebanks within family) to generate 95\% confidence intervals. Benjamini-Hochberg FDR correction applied across all families tested.

\textbf{Results}: Most families show wide, overlapping confidence intervals. NW-Caucasian shows a clear positive residual (point est. 3.62, CI $[3.15, 4.09]$), and Unclassified (polyglot collection) and Indo-Aryan show substantial positive residuals. However, only families with $\geq 3$ treebanks in the bootstrap sample have meaningful CIs; singleton families cannot be reliably ranked.

\section{Results}

\subsection{Primary Finding: No Significant Register Effect at Gold-Label Quality}

\begin{figure}[!htbp]
  \centering
  \includegraphics[width=\linewidth,height=0.85\textheight,keepaspectratio]{figures/fig1_v0.pdf}
  \caption{Non-parametric survival curves (1 minus cumulative hazard) for gold-labeled spoken vs. written dependency arcs across English, French, and Slovenian. Curves show the probability that an arc of length \textgreater d has not yet closed by distance d. Spoken (orange) and written (blue) curves largely overlap within each language, indicating no systematic register difference in arc-length distributions at gold-label quality.}
  \label{fig:fig1}
\end{figure}

Kaplan-Meier survival curves for gold-labeled English, French, and Slovenian show substantial overlap between spoken and written hazard profiles within each language (Figure~\ref{fig:fig1}). The primary Cox model on this subset yields a non-significant register coefficient ($\beta=-0.032$, $p=0.366$). This directly contradicts the hypothesis that spoken language exhibits front-loaded hazard; instead, the gold-labeled data show no systematic register difference in arc-length distribution.

\subsection{Secondary Finding: Apparent Register Effect in Full Corpus is Label-Confounded}

In the full 350-treebank corpus with heuristic labels, a statistically significant register effect emerges ($\beta=+0.046$, $p=1.1\times10^{-4}$). However, this effect is fragile. Label-noise sensitivity analysis shows:

\begin{itemize}
\item 0\% label noise: $\beta=0.011$, $p=0.004$ (significant)
\item 5\% label noise: $\beta=0.007$, $p=0.054$ (marginal)
\item 10\% label noise: $\beta=0.013$, $p=0.0009$ (significant)
\item 20\% label noise: $\beta=0.005$, $p=0.157$ (non-significant)
\end{itemize}

At 20\% perturbation---a plausible noise rate for heuristic labels applied to 347 of 350 treebanks---the effect vanishes. This suggests the full-corpus effect is driven by label assignment bias, not genuine register differences.

\subsection{Strong Typological Effect: Word Order}

\begin{figure}[!htbp]
  \centering
  \includegraphics[width=\linewidth,height=0.85\textheight,keepaspectratio]{figures/fig2_v0.pdf}
  \caption{Estimated Cox regression coefficients from the full-corpus model (350 treebanks, 14.56M arcs) for register (heuristic-labeled), word-order typology, and morphological richness. Point estimates and 95\% confidence intervals shown. Register effect is small and label-noise-dependent (orange, significant in full corpus but confounded); word-order effect is large and highly significant (blue, $p=4.9\times10^{-25}$); morphological richness is not significant (red, $p=0.52$).}
  \label{fig:fig2}
\end{figure}

The word-order coefficient ($\beta=-0.028$, $p=4.9\times10^{-25}$) is large and highly significant (Figure~\ref{fig:fig2}). Free-order languages (low fraction of dependents preceding head) exhibit lower hazard, meaning arcs are less likely to close at short distances---they have flatter, lower-peak hazard curves. Fixed-order languages (high fraction preceding) show steeper hazard, with closure concentrated at shorter distances.

Effect size: A one-standard-deviation increase in word-order scale (from fixed to free) corresponds to a hazard ratio of $\exp(-0.028)=0.972$, a 2.8\% decrease in instantaneous closure risk. While the percentage is small, the effect spans an entire typological dimension and is observed across 14.56 million arcs.

Functional vs. lexical stratification: Functional dependencies (articles, case markers) show weaker register effects ($\beta=0.027$, CI $[0.018, 0.036]$, $p=1.6\times10^{-8}$) than lexical dependencies ($\beta=0.122$, CI $[0.115, 0.129]$, $p=2.7\times10^{-257}$), a 4.5$\times$ ratio consistent with Gerdes et al. \citep{Gerdes2026}.

\subsection{Family-Level Structure}

\begin{figure}[!htbp]
  \centering
  \includegraphics[width=\linewidth,height=0.85\textheight,keepaspectratio]{figures/fig3_v0.pdf}
  \caption{Point estimates and 95\% bootstrap confidence intervals for residual Nelson-Aalen cumulative hazard at d=10 across language families. Families are sorted by point estimate. Families with fewer than 2 treebanks (insufficient for bootstrap CI) are omitted. Wide CIs reflect small sample sizes; only families with \textgreater=5 treebanks have narrow CIs. NW-Caucasian and Unclassified show clear positive residuals; Romance and Slavic show negative residuals. Most CIs overlap zero, suggesting family-level heterogeneity is modest after typological covariates are controlled.}
  \label{fig:fig3}
\end{figure}

Language families show substantial heterogeneity in residual hazard after word-order and morphological-richness covariates are controlled (Figure~\ref{fig:fig3}). Bootstrap CIs are wide for most families (singleton or small-sample families), but a few show consistent positive or negative residuals. NW-Caucasian shows notably elevated hazard relative to its typological cluster, while Romance, Slavic, and Indo-Aryan show lower-than-expected hazard.

Interpretation: Family-level deviations suggest language families have distinct grammatical or processing strategies that go beyond word-order typology alone. However, sample-size constraints limit the strength of these claims; replication on larger family-level samples is necessary.

\subsection{Cross-Check Against Futrell et al. and Gerdes et al.}

The hypothesis predicted recovery of Futrell et al. (2015)'s finding that all 37 languages minimize dependency length vs. random baseline. A random-head-permutation null (heads reassigned uniformly within sentence-length bounds) yields mean arc length 8.77 vs. 3.38 observed, a clear and large separation (Nelson-Aalen AUC difference 78.8). This replicates Futrell's directional result: DLM is strong and consistent across our 350-treebank sample.

The functional/lexical split is confirmed: functional dependencies ($\beta=0.027$) show weaker language effects than lexical dependencies ($\beta=0.122$), consistent with Gerdes et al.'s hypothesis that grammar-driven dependencies are universal while processing-driven dependencies are typologically variable.

\section{Discussion}

\subsection{Methodological Contribution: Survival Analysis as a Confound-Resolution Tool}

The primary contribution of this work is methodological: survival-analysis methods provide a principled, built-in solution to the length-mixing confound that has long plagued dependency-length research. By treating sentence-boundary constraints as censoring (not as a regression predictor to normalize away), we eliminate the mechanical confound at its source. This reframing is not novel to dependency data---biostatistics has used it for decades---but its application to synchronic linguistic data is, to our knowledge, unprecedented.

The robustness check partially validates this advantage: Cox coefficients are more stable under sentence-length-composition resampling than pooled-MDD ratios. However, the quantitative advantage (1.3$\times$ variance ratio, not 10--20$\times$) is more modest than originally hypothesized, suggesting the confound's practical impact may be smaller in some regimes than others.

\subsection{The Register Finding: A Cautionary Tale on Label Quality}

Our analysis reveals a stark contrast between gold-labeled and heuristic-labeled registers:
\begin{itemize}
\item \textbf{Gold-labeled subset} ($n=25{,}710$, 3 languages): $\beta=-0.032$, $p=0.366$ (not significant).
\item \textbf{Full-corpus heuristic labels} ($n=14.56$M, 350 treebanks): $\beta=+0.046$, $p=1.1\times10^{-4}$ (significant, but label-noise-dependent).
\end{itemize}

This 146\% discrepancy and label-noise sensitivity are significant findings in themselves. They demonstrate that register effects in dependency-length research are highly sensitive to annotation quality. For future work, we recommend:
\begin{enumerate}
\item Prioritize gold-labeled spoken/written corpora (CHILDES, Rhapsodie, SST, etc.) over heuristic labeling.
\item Explicitly model label uncertainty, rather than treating register as a fixed covariate.
\item Report both gold-labeled and heuristic results, with transparent quality flags.
\end{enumerate}

Our honest finding is that \textbf{spoken language does not show significantly front-loaded dependency hazard at gold-label quality}. The apparent effect in the full corpus is confounded by label assignment bias. This does not invalidate the register hypothesis; rather, it underscores that the hypothesis needs cleaner data to test.

\subsection{Typological Effects: Robust and Large}

The word-order effect ($\beta=-0.028$, $p=4.9\times10^{-25}$) survives all robustness checks and operationalization variants. Free-order languages exhibit flatter, lower-peak hazard, consistent with the idea that morphological marking (case, agreement) permits longer dependencies without real-time ambiguity. This is a genuine typological signal. Effect size, while a 2.8\% hazard decrease per SD, is meaningfully large at the 14.56M-arc scale and aligns with linguistic theory.

\subsection{Family-Level Structure: Tentative and Exploratory}

Language families show residual heterogeneity, but bootstrap CIs are wide for most families due to limited treebank coverage. NW-Caucasian emerges as an outlier, but this is based on a small sample ($n_{\text{treebanks}}=2$). We caution against over-interpreting family rankings without larger, more balanced language-family samples in UD.

\subsection{Limitations}

\begin{enumerate}
\item \textbf{Register labeling}: Only 3 of 350 treebanks have gold-documented spoken/written splits. The primary register analysis is restricted to these 3 languages, limiting generalizability. The full-corpus heuristic-labeled estimate is confounded by label noise.

\item \textbf{Word-order operationalization}: The primary covariate is empirical (fraction of dependents preceding head), not categorical (Grambank/WALS). While this ensures consistency, it differs from typological classifications linguists may prefer. A future sensitivity analysis should compare against Grambank categorical classes.

\item \textbf{Family-level frailty}: We use stratification and post-hoc residual ranking as proxies for family-level frailty, rather than explicit random-effect frailty modeling (which lifelines does not natively support). Bayesian methods (e.g., PyMC) would provide more rigorous family-level inference.

\item \textbf{Functional vs. lexical stratification}: Register effects are larger for lexical dependencies ($\beta=0.122$) than functional ($\beta=0.027$). The functional-dependency hypothesis---that grammar-driven dependencies are universal---is supported, but the analysis does not deeply examine why lexical dependencies show larger register variance. Deeper linguistic modeling (e.g., by argument structure, semantic role) could refine this.

\item \textbf{Sample-size asymmetry}: Spoken arcs are far fewer than written (12,855 vs. 12,855 in gold subset, but 18,846 vs. 67,434 in full corpus across all languages). Small spoken samples in many languages limit power for language-specific register effects.
\end{enumerate}

\section{Conclusion}

We have introduced survival-analysis methods to the study of dependency-length minimization, treating arc length as a right-censored time-to-event outcome. This reframes the documented sentence-length-pooling confound as a design component of the model, rather than a statistical hazard to be normalized away.

Our analysis of 14.56 million arcs across 350 UD treebanks yields three findings:

\begin{enumerate}
\item \textbf{Methodological}: Survival analysis provides a confound-robust framework for quantitative typology. Cox regression coefficients are stable under sentence-length-composition resampling, validating the reframing's utility.

\item \textbf{Register}: On gold-labeled spoken/written pairs, no significant register effect emerges ($\beta=-0.032$, $p=0.366$). The apparent effect in the full corpus ($\beta=+0.046$) is confounded by heuristic register labeling and vanishes under label-noise sensitivity analysis. Future research should prioritize gold-labeled spoken corpora.

\item \textbf{Typology}: Word-order class is a robust, large predictor of hazard-curve shape ($\beta=-0.028$, $p=4.9\times10^{-25}$). Free-order languages exhibit flatter profiles, consistent with theories linking morphological richness to dependency-length tolerance. Language families exhibit residual structure beyond typology, though bootstrap CIs require larger samples for definitive conclusions.
\end{enumerate}

This work opens a new methodological avenue for quantitative typology, demonstrating that survival-analysis tools can be adapted to linguistic problems with hidden censoring structures. It also serves as a cautionary example: apparent large-scale effects can be artifacts of label quality, emphasizing the need for transparent data provenance in linguistic research.

\bibliographystyle{plainnat}
\bibliography{references}

\end{document}
